\begin{theorem}[All marked limiting period groups]\hypertarget{tag-032H}{}\label{push:src:D.3.1}
\StatementRecord{push:src:D.3.1}{proved}{032H}{\SeriesRef{push:src:D.2.1}, \SeriesRef{eq:d-dual-exact}}

Let \(c_n\to c_*\) with \(\operatorname{Im}c_n<H_*\), and let
\(h_n\to0\). Set
\[
 \delta_n=-D_{c_n}(h_n)>0,\qquad
 V_n=\left(\frac{6\operatorname{Re}h_n}{\delta_n},
           \frac{6\operatorname{Im}h_n}{\delta_n},
           \frac{\operatorname{Re}(c_n-c_*)}{\delta_n}\right).
\]
Define the discrete group
\[
 L_a=\{(p,q,r)\in\mathbb R^3:
    r\in\mathbb Z,\ p+r/2\in\mathbb Z,
    \ q+r/2\in\mathbb Z,\ ar\in\mathbb Z\}.
\]
It has rank three for rational \(a\), and is
\(\mathbb Z^2\times\{0\}\) for irrational \(a\).
Every convergent subsequence of the closed lattices
\(\Gamma_n=\Pi_{c_n}(h_n)\mathbb Z^4\subset\mathbb C^2\) has the
following form. There is a primitive subgroup \(K\subset L_a\) and
a homomorphism \(\ell:K\to\mathbb R\) such that
\begin{equation}\hypertarget{tag-01OM}{}\label{eq:d-limit-group}
 \Gamma_\infty=
 \{(\eta,\xi):p\operatorname{Re}\eta+q\operatorname{Im}\eta
       +r\operatorname{Re}\xi+\ell(p,q,r)\operatorname{Im}\xi
       \in\mathbb Z\text{ for every }(p,q,r)\in K\}.
\end{equation}
Here
\[
 K=\{\lambda\in L_a:\lambda\cdot V_n
                      \text{ has a finite limit}\},\qquad
 \ell(\lambda)=\lim_n\lambda\cdot V_n,
\]
and for \(\lambda\notin K\) the absolute value tends to infinity.
Conversely every primitive \(K\subset L_a\) and every such
\(\ell\) occur for sequences with \(\operatorname{Im}c_n<H_*\).
Thus the possible pointed additive quotient limits are real tori of
any dimension from zero to three for rational \(a\), and from zero
to two for irrational \(a\), with the marked groups given explicitly
by \eqref{eq:d-limit-group}.
\end{theorem}


\begin{proof}
\begin{enumerate}[label=\arabic*.,leftmargin=*]
\item We use the fixed complex-linear coordinates of \ref{push:src:D.2.1}. Write
\(h=X+iY\), \(\tau=R+iT\), and
\(z=B(h)+c-c_*\). The real period matrix, with the basis of \ref{push:src:D.2.1},
is
\[
 A=\begin{pmatrix}
 1&R&0&6X\\
 0&T&0&6Y\\
 1/2&1/2+X&1&a+\operatorname{Re}z\\
 0&Y&0&\operatorname{Im}z
 \end{pmatrix},\qquad
 \det A=T D_c.
\]
For an integral dual label \((m_w,m_u,m_d,m_g)\), put
\(r=m_d\), \(p=m_w-r/2\), \(q=m_u-r/2\).
The real-linear covector \(A^{-t}(m_w,m_u,m_d,m_g)^t\) is
\begin{equation}\hypertarget{tag-01OL}{}\label{eq:d-dual-exact}
 \left(p,\frac{q-pR-rX-tY}{T},r,t\right),\quad
 t=\frac{m_g-ra-6pX-6Y(q-pR-rX)/T-r\operatorname{Re}z}{D_c}.
\end{equation}
Multiplication by \(A^t\) proves the formula.

\item A bounded sequence of dual vectors has bounded integral labels,
since \(A_n\) is bounded. After subselection its label is constant.
Its limit must satisfy \(m_g=ra\), because
\(A_0^t\) has that relation. The remaining labels are exactly
\(L_a\). Conversely consider a fixed label from \(L_a\).
Lemma~\ref{push:src:D.2.1} gives
\[
 T=1+O(|h|^2),\quad R=O(|h|^2),\quad
 \operatorname{Re}B=6XY+O(|h|^4),
\]
\[
 \delta_n=-\operatorname{Im}(c_n-c_*)
                       +3|h_n|^2+O(|h_n|^4).
\]
In particular \(\delta_n\ge C|h_n|^2\) for a fixed \(C>0\)
after shrinking the disc. The numerator in
\eqref{eq:d-dual-exact}, with \(m_g=ra\), is
\[
 -6pX-6qY-r\operatorname{Re}(c-c_*)
       +O(|h|^3)(|p|+|q|+|r|).
\]
The cancellation of the quadratic real term is essential. Division
by \(D_c=-\delta_n\) makes the error \(O(|h_n|)\). Thus
\[
 t_n=(p,q,r)\cdot V_n+o(1).
\]
Whenever this is bounded, the other three coefficients converge to
\((p,q,r)\). This proves that the dual limit, if it exists, is the
discrete graph
\(\{(p,q,r,\ell(p,q,r)):(p,q,r)\in K\}\).
There can be no further elements from varying labels, because a
bounded sequence of labels lies in a finite set. A convergent dual
group cannot contain two different limits with the same label: their
difference would have label zero, whose sole preimage under every
\(A_n^t\) is zero. Nor can a surviving label escape along a
subsequence, since an element of a Chabauty limit must be approximable
along the whole convergent sequence. Consequently every label either
has the indicated finite limit or escapes to infinity. Addition and
integer division show that \(K\) is a subgroup and is primitive.

\item By \cite[Theorem~3.1, p.~2017]{Cornulier2011}, the map from
closed subgroups of a finite-dimensional real vector space to their
Pontryagin annihilators is a homeomorphism. Applying it to the dual graph gives exactly
\eqref{eq:d-limit-group}. This theorem's hypotheses are met by the
real vector space underlying \(\mathbb C^2\) and by the closed full
lattices \(\Gamma_n\); no nonclosed group is used as a Chabauty
point.

\item To prove realizability, extend \(\ell\) linearly to a covector
represented by a vector \(v_0\in\mathbb R^3\). Since \(K\) is
primitive, \(L_a\cap\operatorname{span}_{\mathbb R}K=K\).
Choose \(v\) orthogonal to that span and not orthogonal to any member
of \(L_a\setminus K\); the excluded proper hyperplanes are
countably many, and cannot cover the orthogonal space. If
\(K=L_a\), take \(v=0\). Put \(V_n=nv+v_0=(A_n',B_n',C_n')\).
Choose positive \(d_n\to0\) with
\(d_n(1+\|V_n\|^2)\to0\), and define
\[
 h_n=d_n(A_n'+iB_n')/6,\quad x_n=d_nC_n',\quad
 \epsilon_n=d_n+D_{c_*}(h_n)>0,
 \qquad c_n=c_*+x_n-i\epsilon_n.
\]
The inequality follows from \(D_{c_*}(h)=-3|h|^2+O(|h|^4)\).
These choices give \(-D_{c_n}(h_n)=d_n\) exactly, so their vectors
are the prescribed \(V_n\). On \(K\) the scalar products converge
to \(\ell\), and off \(K\) they escape. This constructs every
claimed limit inside the admissible parameter domain.

\item Finally the dual graph is a discrete free group of rank
\(\operatorname{rank}K\). Its annihilator quotient is a compact
connected real torus of that dimension. In the fixed Euclidean metric,
\(d_n([z],[z'])=\operatorname{dist}(z-z',\Gamma_n)\).
Chabauty convergence gives local uniform convergence of these
functions: minimizers relevant on a bounded set can be restricted to
a bounded set, and the two defining Chabauty approximation properties
then give both inequalities. This identifies the stated pointed
quotient limits. Their group structures and quotient metrics are the ones displayed above.
\end{enumerate}
\end{proof}